% body.tex —— 导学件答案册-v1（组装器 组装body.py 产出，可重跑；勿手改——改口径请先改组装器）
% 内容＝去答案0911 移出块全量（导学件正文判分值 0 在场，本册值在场）；题号与导学件同号。
% 页码＝独立册起页 1（样张全书连续 P75 系练习册装订档，不适用，登记简报）。
\setcounter{page}{1}
% ---------- 册首行（沿样张拍板⑨：册名「参考答案」，E2 节标题件承载） ----------
\par\nointerlineskip
{\centering\fontsize{16.88pt}{22pt}\selectfont\hejie 参考答案\par}
\vspace{2.4mm}
\zhangtitle{第一章\quad 空间向量与立体几何}
\jietitle{1.1.1\quad 空间向量及其运算(导学件答案)}
\begin{multicols}{2}
\raggedcolumns
\emergencystretch=1em
% ---------- 件：CJK 题号条目（样张 \ansitem 承形——「例1/变式1」非数字号，题号走 \heihao 中文族；v10-B 层级缩进＝续行/后续段收一档 \qpind） ----------
\newcommand{\dansitem}[2]{\par\glueguard{1}\addvspace{4pt}%
  {\hangindent\qpind\hangafter=1\setlength{\parindent}{\qpind}%
  \noindent
  \makebox[0pt][l]{{\fontsize{11.4pt}{13pt}\selectfont\heihao #1}}%
  \hspace*{\anshang}{\anlabel [答案]}\hspace{0.5em}#2\par}}
\newcommand{\txline}[1]{\par\noindent\hspace*{\qpind}%
  {\kaishu\fontsize{10.09pt}{12pt}\selectfont 题型：#1}\par}
\newcommand{\txsummaryhead}[1]{\par\glueguard{1}\addvspace{3pt}\noindent%
  {\heibf\fontsize{10.09pt}{13pt}\selectfont [题型总结]\hspace{0.5em}#1}\par\nopagebreak}
\newcommand{\txsumitem}[2]{\par\noindent\hspace*{\qpind}%
  {\kaishu\fontsize{10.09pt}{12pt}\selectfont #1：#2}\par}
\qufen{一、课前预习}{知识点填空与判断答案——序与导学件同}

\kdhead{知识点一　空间向量的概念}{填空11空／判断2题——序与导学件同}
\ansline{答案}{1．大小　方向

2．长度

3．零　0　0　1　相反　互相平行　重合　模相等

(1)×　(2)√}
\ansline{解析}{(1)模相等方向未必相同，还可能相反或斜交，两向量未必相等，故×.

(2)相等向量方向相同，所在直线平行或重合，符合共线向量(平行向量)定义，故√.}
\kdhead{知识点二　空间向量的线性运算}{填空5空／判断2题——序与导学件同}
\ansline{答案}{1．三角形　平行四边形　相同　相反

2．实数

(1)×　(2)×}
\ansline{解析}{(1)平行传递需\(\overrightarrow{b}\neq\overrightarrow{0}\)；\(\overrightarrow{b}=\overrightarrow{0}\)时零向量与任意向量平行，传递性失效，故×.

(2)缺\(\overrightarrow{b}\neq\overrightarrow{0}\)前提；且\(\overrightarrow{a}=\overrightarrow{0}\)时\(\lambda\)不唯一，故×.}
\kdhead{知识点三　空间向量的夹角、数量积与共面}{填空4空／判断2题——序与导学件同}
\ansline{答案}{1．[0°,180°]　互相垂直

2．零向量

4．有序实数对

(1)√　(2)×}
\ansline{解析}{(1)\(\overrightarrow{p}\)是\(\overrightarrow{a}\)、\(\overrightarrow{b}\)的线性组合，由共面向量定理知三向量共面，故√.

(2)数量积大于0仅保证\(\cos\theta>0\)，夹角\(\theta\in[0^\circ,90^\circ)\)；\(\theta=0^\circ\)(同向)时数量积也为正而非锐角，故×.}
\setlength{\anshang}{12.2mm}
\qufen{二、课中探究}{九探究点×例1／变式1——题号与导学件同号}

% ---------- 探究点一　空间向量的概念辨析 ----------
\qufen{探究点一　空间向量的概念辨析}{例1／变式1}
\dansitem{例1}{B}
\txline{空间向量的概念辨析}
\ansline{分析}{根据零向量、单位向量、向量的模、相等向量、相反向量等概念逐一分析即可得出答案.}
\ansline{详解}{解：若\(| \overrightarrow{a} | = | \overrightarrow{b} |\)，则它们的方向相同时是相等向量，方向相反时是相反向量，还有可能方向既不相同，也不相反，A错；

若\(\overrightarrow{a},\overrightarrow{b}\)为相反向量，则它们的和为零向量，B对；

零向量的方向是任意的，C错；

两个单位向量只是模都为1，方向不一定相同，D错.故选：B.}
\dansitem{变式1}{D}
\txline{空间向量的概念辨析}
\ansline{解析}{A规定成立；B模为0即零向量；C相等必等模；D平行不需等模(反例\(\overrightarrow{a}=2\overrightarrow{b}\))，故选D.}
\txsummaryhead{空间向量的概念辨析}
\txsumitem{识别信号}{题干围绕模、方向、零向量、单位向量、相等向量与相反向量罗列说法要求辨误时，判归本题型.}
\txsumitem{通法步骤}{对照定义逐条核验，存疑条目用反例否定.}
\txsumitem{易错点}{排除错误项定答案，忌凭直觉误选.}

% ---------- 探究点二　空间向量的线性运算 ----------
\qufen{探究点二　空间向量的线性运算}{例1／变式1}
\dansitem{例1}{\(\overrightarrow{b} - \overrightarrow{a} - \overrightarrow{c}\)}
\txline{空间向量的线性运算}
\ansline{分析}{连接\(CA_{1},\)根据直三棱柱的结构特征及空间向量减法的几何意义可得\(\overrightarrow{A_{1}B} = \overrightarrow{CB} - \overrightarrow{CA} - \overrightarrow{CC_{1}}\)，结合已知即可求表达式.}
\ansline{详解}{连接\(CA_{1},\)则\nobreak\(\overrightarrow{A_{1}B} = \overrightarrow{CB} - \overrightarrow{CA_{1}} = \overrightarrow{CB} - \overrightarrow{CA} - \overrightarrow{CC_{1}} =\overrightarrow{b} - \overrightarrow{a} - \overrightarrow{c}\).

故答案为：\(\overrightarrow{b} - \overrightarrow{a} - \overrightarrow{c}\)}
\ansfig{\includegraphics[width=28.800mm]{C:/提示词/工作区/字替对照-0909/variantF/media/media/image1.png}}
\dansitem{变式1}{\(\overrightarrow{a}-\overrightarrow{b}-\overrightarrow{c}\)}
\txline{空间向量的线性运算}
\ansline{解析}{\(\overrightarrow{D_1B}=\overrightarrow{D_1A_1}+\overrightarrow{A_1A}+\overrightarrow{AB}=-\overrightarrow{b}-\overrightarrow{c}+\overrightarrow{a}\).}
\txsummaryhead{空间向量的线性运算}
\txsumitem{识别信号}{在几何体中给出若干已知向量、要求表示指定向量时，判归本题型.}
\txsumitem{通法步骤}{为目标向量选一条首尾相接的路径写出加减链，再把链中各段用相等或共线的已知向量替换，最后合并化简得表达式.}
\txsumitem{易错点}{复查每段方向与起讫点，防抄错向量.}

% ---------- 探究点三　用基底表示向量 ----------
\qufen{探究点三　用基底表示向量}{例1／变式1}
\dansitem{例1}{\(1\)；\(\frac{1}{4}\)}
\txline{用基底表示向量}
\ansline{分析}{结合空间向量的线性运算列方程，由此求得\(x,y\)的值.}
\ansline{详解}{\(\overrightarrow{AE} = \overrightarrow{AA_{1}} + \overrightarrow{A_{1}E} = \overrightarrow{AA_{1}} + \frac{1}{4}\overrightarrow{A_{1}C_{1}} = \overrightarrow{AA_{1}} + \frac{1}{4}( \overrightarrow{AB} + \overrightarrow{AD} )\)，所以\(x = 1,y = \frac{1}{4}\).故答案为：\(1\)；\(\frac{1}{4}\)}
\dansitem{变式1}{\(x=1\)，\(y=\frac{1}{2}\)}
\txline{用基底表示向量}
\ansline{解析}{\(\overrightarrow{AE}=\overrightarrow{AA_1}+\overrightarrow{A_1E}=\overrightarrow{AA_1}+\frac{1}{2}(\overrightarrow{AB}+\overrightarrow{AD})\)，故\(x=1\)，\(y=\frac{1}{2}\).}
\txsummaryhead{用基底表示向量}
\txsumitem{识别信号}{指定三个不共面的向量作基底、要求把目标向量写成基底向量的线性组合并求系数时，判归本题型.}
\txsumitem{通法步骤}{先把目标向量沿几何路径分解为基底向量的和，再把各段按比例关系代入基底.}
\txsumitem{易错点}{比较等式两边系数求出待定值，防漏写中间段.}

% ---------- 探究点四　共面向量的判定 ----------
\qufen{探究点四　共面向量的判定}{例1／变式1}
\dansitem{例1}{C}
\txline{共面向量的判定}
\ansline{分析}{利用空间向量共面定理逐一判断即可.}
\ansline{详解}{若\(\overrightarrow{a}\)，\(\overrightarrow{b}\)共线，由\(\overrightarrow{p}\)=\emph{x}\(\overrightarrow{a}\)+\emph{y}\(\overrightarrow{b}\)知\(\overrightarrow{p}\)一定与\(\overrightarrow{a}\)，\(\overrightarrow{b}\)共面，

若\(\overrightarrow{a}\)，\(\overrightarrow{b}\)不共线，则满足共面定理，\(\overrightarrow{p}\)与\(\overrightarrow{a}\)，\(\overrightarrow{b}\)共面，①对；

同理③对；若\(\overrightarrow{p}\)与\(\overrightarrow{a}\)，\(\overrightarrow{b}\)共面，且\(\overrightarrow{a}\)，\(\overrightarrow{b}\)共线，则不一定有\(\overrightarrow{p}\)=\emph{x}\(\overrightarrow{a}\)+\emph{y}\(\overrightarrow{b}\)，故②不对；同理④不对，故选：C.}
\dansitem{变式1}{共面(\(\overrightarrow{MP}\)可由\(\overrightarrow{MA}\)，\(\overrightarrow{MB}\)线性表示，故\(\overrightarrow{MP}\)，\(\overrightarrow{MA}\)，\(\overrightarrow{MB}\)共面，即\(P\)在平面\(MAB\)内)}
\txline{共面向量的判定}
\ansline{解析}{\(\overrightarrow{MP}=3\overrightarrow{MA}-2\overrightarrow{MB}\)为\(\overrightarrow{MA}\)、\(\overrightarrow{MB}\)的线性组合，故三向量共面，\(P\)在平面\(MAB\)内.}
\txsummaryhead{共面向量的判定}
\txsumitem{识别信号}{题干给出p=xa+yb或MP=xMA+yMB一类线性表示与“共面”结论的互推说法判断时，判归本题型.}
\txsumitem{通法步骤}{核对充要条件的前提(两向量不共线或三点不共线)是否齐备，再对缺前提的说法举共线反例排除.}
\txsumitem{易错点}{最后汇总正确序号定选项.}

% ---------- 探究点五　数量积的概念与运算律 ----------
\qufen{探究点五　数量积的概念与运算律}{例1／变式1}
\dansitem{例1}{AD}
\txline{数量积的概念与运算律}
\ansline{分析}{根据空间向量数量积的定义与运算律一一判断即可；}
\ansline{详解}{解：对于A：\({\overrightarrow{a}}^{2} = \overrightarrow{a} \cdot \overrightarrow{a} = | \overrightarrow{a} | \cdot | \overrightarrow{a} |\cos 0 = | \overrightarrow{a} |^{2}\)，故A正确；

对于B：因为向量不能做除法，即\(\frac{\overrightarrow{b}}{\overrightarrow{a}}\)无意义，故B错误；

对于C：\(( \overrightarrow{a} \cdot \overrightarrow{b} )^{2} = ( | \overrightarrow{a} | \cdot | \overrightarrow{b} |\cos\langle \overrightarrow{a},\overrightarrow{b} \rangle )^{2} = | \overrightarrow{a} |^{2} \cdot | \overrightarrow{b} |^{2}\cos^{2}\langle \overrightarrow{a},\overrightarrow{b} \rangle\)，故C错误；

对于D：\(( \overrightarrow{a} - \overrightarrow{b} )^{2} = ( \overrightarrow{a} - \overrightarrow{b} ) \cdot ( \overrightarrow{a} - \overrightarrow{b} ) = {\overrightarrow{a}}^{2} - 2\overrightarrow{a} \cdot \overrightarrow{b} + {\overrightarrow{b}}^{2}\)，故D正确；故选：AD}
\dansitem{变式1}{\(-11\)}
\txline{数量积的概念与运算律}
\ansline{解析}{\(\overrightarrow{a}\cdot\overrightarrow{b}=2\times3\times\cos60^\circ=3\)；原式\(=\overrightarrow{a}^2+\overrightarrow{a}\cdot\overrightarrow{b}-2\overrightarrow{b}^2=4+3-18=-11\).}
\txsummaryhead{数量积的概念与运算律}
\txsumitem{识别信号}{题干列出数量积相关的若干等式(平方、商式、乘积展开)要求判断正误时，判归本题型.}
\txsumitem{通法步骤}{按定义a·b=|a||b|cosθ逐式展开核验，展开后移项验算.}
\txsumitem{易错点}{警惕向量无除法、数量积无结合律等误用.}

% ---------- 探究点六　数量积求夹角与投影 ----------
\qufen{探究点六　数量积求夹角与投影}{例1／变式1}
\dansitem{例1}{\(60{^\circ}\)}
\txline{数量积求夹角与投影}
\ansline{分析}{方法1：结合图形，平移向量，利用空间向量的夹角定义来求，但要注意向量夹角的范围；

方法2：先求\(\overrightarrow{BC_{1}} \cdot \overrightarrow{AC}\)，再利用公式\(\cos\langle \overrightarrow{BC_{1}},\overrightarrow{AC} \rangle = \frac{\overrightarrow{BC_{1}} \cdot \overrightarrow{AC}}{| \overrightarrow{BC_{1}} | \cdot | \overrightarrow{AC} |}\)求\(\cos\langle \overrightarrow{BC_{1}},\overrightarrow{AC} \rangle\)，最后确定\(\langle \overrightarrow{BC_{1}},\overrightarrow{AC} \rangle\)即可.}
\ansline{详解}{解：方法1：因为\(\overrightarrow{AD_{1}} = \overrightarrow{BC_{1}}\)，所以\(\angle CAD_{1}\)的大小就等于\(\langle \overrightarrow{BC_{1}},\overrightarrow{AC} \rangle\)

因为△\(CAD_{1}\)为等边三角形，所以\(\angle CAD_{1} = 60^{{^\circ}}\)，所以\(\overrightarrow{BC_{1}}\)与\(\overrightarrow{AC}\)的夹角的大小为\(60{^\circ}\).

方法2．设正方体的棱长为1，\(\overrightarrow{BC_{1}} \cdot \overrightarrow{AC} = ( \overrightarrow{BC} + \overrightarrow{CC_{1}} ) \cdot ( \overrightarrow{AB} + \overrightarrow{BC} ) = ( \overrightarrow{AD} + \overrightarrow{AA_{1}} ) \cdot ( \overrightarrow{AB} + \overrightarrow{AD} )\)\allowbreak\(= \overrightarrow{AD} \cdot \overrightarrow{AB} + | \overrightarrow{AD} |^{2} + \overrightarrow{AA_{1}} \cdot \overrightarrow{AB} + \overrightarrow{AA_{1}} \cdot \overrightarrow{AD} = 0 + | \overrightarrow{AD} |^{2} + 0 + 0 = | \overrightarrow{AD} |^{2} = 1\)又因为\(| \overrightarrow{BC_{1}} | = \sqrt{2},| \overrightarrow{AC} | = \sqrt{2}\)，所以\(\cos\langle \overrightarrow{BC_{1}},\overrightarrow{AC} \rangle = \frac{\overrightarrow{BC_{1}} \cdot \overrightarrow{AC}}{| \overrightarrow{BC_{1}} | \cdot | \overrightarrow{AC} |} = \frac{1}{\sqrt{2} \times \sqrt{2}} = \frac{1}{2}\)，
因为\(\langle \overrightarrow{BC_{1}},\overrightarrow{AC} \rangle \in \lbrack 0,\pi\rbrack\)，所以\(\overrightarrow{BC_{1}}\)与\(\overrightarrow{AC}\)的夹角的大小为\(60{^\circ}\).}
\dansitem{变式1}{\(\frac{\sqrt{10}}{5}\)}
\txline{数量积求夹角与投影}
\ansline{解析}{\(\overrightarrow{AB_1}\cdot\overrightarrow{AC}=(\overrightarrow{AB}+\overrightarrow{BB_1})\cdot(\overrightarrow{AB}+\overrightarrow{BC})=|\overrightarrow{AB}|^2=4\)，\(|\overrightarrow{AB_1}|=\sqrt{5}\)，\(|\overrightarrow{AC}|=2\sqrt{2}\)，故\(\cos\langle\overrightarrow{AB_1},\overrightarrow{AC}\rangle=\frac{4}{\sqrt{5}\times2\sqrt{2}}=\frac{\sqrt{10}}{5}\).}
\txsummaryhead{数量积求夹角与投影}
\txsumitem{识别信号}{题干在几何体中求两向量的夹角、模长或投影而不建坐标系时，判归本题型.}
\txsumitem{通法步骤}{先平移向量使起点共点或选不共面的三个向量作基底，再把相关向量用基底表示并算出数量积与模长.}
\txsumitem{易错点}{最后代入夹角公式并按向量夹角范围[0,π]取角.}

% ---------- 探究点七　数量积条件求参 ----------
\qufen{探究点七　数量积条件求参}{例1／变式1}
\dansitem{例1}{\(( - 1 - \sqrt{3}, - 1 + \sqrt{3} )\)}
\txline{数量积条件求参}
\ansline{分析}{根据\(( \overrightarrow{a} + \lambda\overrightarrow{b} ) \cdot ( \lambda\overrightarrow{a} - 2\overrightarrow{b} ) < 0\)，得出方程求得\(- 1 - \sqrt{3} < \lambda < - 1 + \sqrt{3}\)，若向量\(\overrightarrow{a} + \lambda\overrightarrow{b}\)与\(\lambda\overrightarrow{a} - 2\overrightarrow{b}\)反向共线时，则存在实数\(k < 0\)使得\(\overrightarrow{a} + \lambda\overrightarrow{b} = k(\lambda\overrightarrow{a} - 2\overrightarrow{b})\)成立，得到方程组\(\Bigg\{ \begin{array}{r}
k\lambda = 1 \\
\lambda = - 2k
\end{array} \Bigg.\)无解，即可得到答案.}
\ansline{详解}{由向量\(\overrightarrow{a},\overrightarrow{b}\)，且\(| \overrightarrow{a} | = 2\)，\(| \overrightarrow{b} | = 1\)，\(\langle \overrightarrow{a},\overrightarrow{b} \rangle = 60^{\circ}\)，可得\(\overrightarrow{a} \cdot \overrightarrow{b} = 2 \times 1 \times \cos 60^{\circ} = 1\)，

因为向量\(\overrightarrow{a} + \lambda\overrightarrow{b}\)与\(\lambda\overrightarrow{a} - 2\overrightarrow{b}\)的夹角为钝角，可得\(\cos\langle \overrightarrow{a} + \lambda\overrightarrow{b},\lambda\overrightarrow{a} - 2\overrightarrow{b} \rangle < 0\)且\(\cos\langle \overrightarrow{a} + \lambda\overrightarrow{b},\lambda\overrightarrow{a} - 2\overrightarrow{b} \rangle \neq - 1\)，由\(( \overrightarrow{a} + \lambda\overrightarrow{b} ) \cdot ( \lambda\overrightarrow{a} - 2\overrightarrow{b} ) < 0\)，即\(\lambda{\overrightarrow{a}}^{2}\text{+(}\lambda^{2} - 2)\overrightarrow{a} \cdot \overrightarrow{b} - 2\lambda{\overrightarrow{b}}^{2} < 0\)，所以\(\lambda^{2} + 2\lambda - 2 < 0\)，解得\(- 1 - \sqrt{3} < \lambda < - 1 + \sqrt{3}\)，

若向量\(\overrightarrow{a} + \lambda\overrightarrow{b}\)与\(\lambda\overrightarrow{a} - 2\overrightarrow{b}\)反向共线时，存在实数\(k < 0\)，使得\(\overrightarrow{a} + \lambda\overrightarrow{b} = k(\lambda\overrightarrow{a} - 2\overrightarrow{b})\)成立，可得\(\Bigg\{ \begin{array}{r}
k\lambda = 1 \\
\lambda = - 2k
\end{array} \Bigg.\)，此时方程组无解，所以不存在实数\(k\)使得向量\(\overrightarrow{a} + \lambda\overrightarrow{b}\)与\(\lambda\overrightarrow{a} - 2\overrightarrow{b}\)反向共线，所以实数\(\lambda\)的取值范围是\(( - 1 - \sqrt{3}, - 1 + \sqrt{3} )\).}
\dansitem{变式1}{\(- \frac{3}{2}\)}
\txline{数量积条件求参}
\ansline{解析}{\(\overrightarrow{m}\cdot\overrightarrow{n}=(\overrightarrow{a}+\overrightarrow{b})\cdot(\overrightarrow{a}+\lambda\overrightarrow{b})=|\overrightarrow{a}|^2+\lambda|\overrightarrow{b}|^2+(1+\lambda)\overrightarrow{a}\cdot\overrightarrow{b}=18+16\lambda-12(1+\lambda)=6+4\lambda=0\)，故\(\lambda=-\frac{3}{2}\).}
\txsummaryhead{数量积条件求参}
\txsumitem{识别信号}{题干给出两向量的模与夹角、并以夹角为锐角/钝角或两向量垂直为条件求参数时，判归本题型.}
\txsumitem{通法步骤}{第一步用数量积把条件翻译成不等式或方程(锐角cos>0、钝角cos<0、垂直cos=0).}
\txsumitem{易错点}{再求解并剔除使两向量共线的参数，最后写出参数值或取值范围.}

% ---------- 探究点八　数量积求距离(展开法) ----------
\qufen{探究点八　数量积求距离(展开法)}{例1／变式1}
\dansitem{例1}{D}
\txline{数量积求距离(展开法)}
\ansline{分析}{由\(\overrightarrow{DB} = \overrightarrow{ED} + \overrightarrow{FE} + \overrightarrow{BF}\)，利用数量积运算性质展开即可得到答案}
\ansline{详解}{\(\because\overrightarrow{BD} = \overrightarrow{ED} + \overrightarrow{FE} + \overrightarrow{BF}\),\(\therefore| \overrightarrow{BD} |^{2} = | \overrightarrow{BF} |^{2} + | \overrightarrow{FE} |^{2} + | \overrightarrow{ED} |^{2} + 2\overrightarrow{BF} \bullet \overrightarrow{FE} + 2\overrightarrow{FE} \bullet \overrightarrow{ED} + 2\overrightarrow{BF} \bullet \overrightarrow{ED} = 1 + 1 + 1 - \sqrt{2}\)故\(| \overrightarrow{BD} | = \sqrt{3 - \sqrt{2}}\)故选\(D\)}
\ansline{点睛}{本题是要求空间两点之间的距离，运用空间向量将其表示，然后计算得到结果，较为基础.}
\dansitem{变式1}{\(\sqrt{2}\)}
\txline{数量积求距离(展开法)}
\ansline{解析}{\(|\overrightarrow{BD}|^2=1+1+1+2\overrightarrow{BF}\cdot\overrightarrow{ED}=3+2\cos120^\circ=2\)，故\(|\overrightarrow{BD}|=\sqrt{2}\).}
\txsummaryhead{数量积求距离(展开法)}
\txsumitem{识别信号}{题干在含二面角的拼接图形中求两点距离时，判归本题型.}
\txsumitem{通法步骤}{第一步把待求线段对应的向量写成已知棱向量的首尾相接和链，再对模长平方作数量积展开，逐项代入模长与夹角(端点处两条垂线段的夹角由二面角决定，注意取45°还是135°).}
\txsumitem{易错点}{最后开方得距离.}

% ---------- 探究点九　数量积求距离(折叠矩形) ----------
\qufen{探究点九　数量积求距离(折叠矩形)}{例1／变式1}
\dansitem{例1}{A}
\txline{数量积求距离(折叠矩形)}
\ansline{分析}{过点\(B\)，\(D\)分别向\(AC\)作垂线，垂足分别为\(M\)，\(N\)，由\(\overrightarrow{BD} = \overrightarrow{BM} + \overrightarrow{MN} + \overrightarrow{ND}\)，平方后结合长度和垂直关系可得解.}
\ansline{详解}{过点\(B\)，\(D\)分别向\(AC\)作垂线，垂足分别为\(M\)，\(N\)，则可得\(AM = \frac{1}{2}\)，\(BM = \frac{\sqrt{3}}{2}\)，\(CN = \frac{1}{2}\)，\(DN = \frac{\sqrt{3}}{2}\)，\(MN = 1\).由于\(\overrightarrow{BD} = \overrightarrow{BM} + \overrightarrow{MN} + \overrightarrow{ND}\)，所以\(|\overrightarrow{BD}|^{2} = {(\overrightarrow{BM} + \overrightarrow{MN} + \overrightarrow{ND})}^{2}\)\allowbreak\(= |\overrightarrow{BM}|^{2} + |\overrightarrow{MN}|^{2} + |\overrightarrow{ND}|^{2} + 2(\overrightarrow{BM} \cdot \overrightarrow{MN} + \overrightarrow{MN} \cdot \overrightarrow{ND} + \overrightarrow{BM} \cdot \overrightarrow{ND})\)\allowbreak\(= ( \frac{\sqrt{3}}{2} )^{2} + 1^{2} + ( \frac{\sqrt{3}}{2} )^{2} + 2 \times (0 + 0 + 0) = \frac{5}{2}\)，所以\(|\overrightarrow{BD}| = \frac{\sqrt{10}}{2}\).故选:\emph{A}.}
\ansfig{\includegraphics[width=42.500mm]{C:/提示词/工作区/字替对照-0909/variantF/media/media/image5.png}}
\ansline{点睛}{本题主要考查了空间向量数量积的应用，属于基础题.}
\dansitem{变式1}{\(\frac{\sqrt{13}}{2}\)}
\txline{数量积求距离(折叠矩形)}
\ansline{解析}{\(BM=DN=\frac{\sqrt{3}}{2}\)，\(MN=1\)；两垂线向量夹角随二面角为\(60^\circ\)，\(\overrightarrow{BM}\cdot\overrightarrow{ND}=\frac{3}{4}\cos60^\circ=\frac{3}{8}\)；\(|\overrightarrow{BD}|^2=\frac{3}{4}+1+\frac{3}{4}+2\times\frac{3}{8}=\frac{13}{4}\)，故\(|\overrightarrow{BD}|=\frac{\sqrt{13}}{2}\).}
\txsummaryhead{数量积求距离(折叠矩形)}
\txsumitem{识别信号}{题干为平面图形沿一条直线折起且折后两面垂直、求两点间距离时，判归本题型.}
\txsumitem{通法步骤}{第一步在折后图形中过两点对折痕作垂线定出垂足，再把目标向量写成两段垂线与折痕段的和链并平方展开.}
\txsumitem{易错点}{最后利用面面垂直带来的向量垂直逐项化简、开方得距离.}

\setlength{\anshang}{5.8mm}
\qufen{三、课堂评价}{本组共5题——题号与导学件同号}

\ansitem{1}{3}
\ansline{解析}{\(\overrightarrow{a}\cdot\overrightarrow{b}=|\overrightarrow{a}||\overrightarrow{b}|\cos60^\circ=2\times3\times0.5=3\).}
\ansitem{2}{B}
\ansline{解析}{\(\overrightarrow{b}=\lambda\overrightarrow{a}\Rightarrow4=2\lambda\Rightarrow\lambda=2\)，\(k=2\)，选B.}
\ansitem{3}{C}
\ansline{解析}{非零向量\(\overrightarrow{a}\cdot\overrightarrow{b}=0\Leftrightarrow\overrightarrow{a}\perp\overrightarrow{b}\)，充要条件，选C.}
\ansitem{4}{D}
\ansline{解析}{取\(AC\)中点\(G\)，\(\overrightarrow{EF}=\overrightarrow{EG}+\overrightarrow{GF}=\frac{1}{2}\overrightarrow{BC}+\frac{1}{2}\overrightarrow{AD}\)，选D.}
\ansitem{5}{\(\sqrt{13}\)}
\ansline{解析}{\(\overrightarrow{a}\cdot\overrightarrow{b}=4\times3\times\cos120^\circ=-6\)，\(|\overrightarrow{a}+\overrightarrow{b}|^2=16+9-12=13\)，故\(\sqrt{13}\).}

\end{multicols}
